\documentclass[a4paper,11pt]{article}
\usepackage[T2A]{fontenc}
\usepackage[utf8]{inputenc}
\usepackage[russian]{babel}
\usepackage{amsmath,amssymb,mathtools}
\usepackage{cmbright}
\usepackage[a4paper,left=22mm,right=21mm,top=20mm,bottom=20mm]{geometry}
\usepackage{microtype,booktabs,tabularx,array,enumitem,parskip}
\usepackage{tikz,xcolor,fancyhdr,hyperref}
\usetikzlibrary{calc}
\definecolor{accent}{HTML}{087C7C}
\definecolor{ink}{HTML}{23333A}
\definecolor{muted}{HTML}{667780}
\definecolor{warn}{HTML}{A43236}
\definecolor{keyblue}{HTML}{245C92}
\hypersetup{colorlinks=true,linkcolor=accent,pdfauthor={Лёвин Артём Александрович},pdftitle={Ревью домашней работы — Ксения Клыкова, 31 октября 2026 г.}}
\setlength{\parindent}{0pt}
\setlength{\parskip}{3pt}
\setlength{\headheight}{14pt}
\renewcommand{\arraystretch}{1.3}
\pagestyle{fancy}
\fancyhf{}
\fancyhead[L]{\small\textcolor{muted}{МАТЕМАТИКА \quad / \quad ТЕОРИЯ ЧИСЕЛ}}
\fancyhead[R]{\small\textcolor{muted}{31.10.2026}}
\fancyfoot[C]{\footnotesize\textcolor{muted}{Лёвин Артём Александрович эксклюзивно для Ксении Клыковой}}
\fancyfoot[R]{\footnotesize\thepage}
\renewcommand{\headrulewidth}{0pt}
\renewcommand{\footrulewidth}{0pt}
\newcommand{\subhead}[1]{\par\vspace{3pt}{\bfseries\textcolor{ink}{#1}}\par\vspace{1pt}}
\newcommand{\errorlabel}{\textcolor{warn}{\textbf{Ошибка:}}}
\newcommand{\keylabel}{\textcolor{keyblue}{\textbf{Ключевая идея:}}}
\newcommand{\finallabel}{\textcolor{accent}{\textbf{Итог:}}}
\newcommand{\tasktitle}[2]{\phantomsection\label{task:#1}\section*{Задание № #1}\textcolor{muted}{#2}\par\vspace{4pt}}
\raggedbottom
\begin{document}
\pagenumbering{roman}
\begin{titlepage}
\thispagestyle{empty}
\centering
\vspace*{23mm}
{\footnotesize\bfseries\textcolor{accent}{МАТЕМАТИКА \; / \; РАБОТА С ДЕЛИТЕЛЯМИ}}\par
\vspace{23mm}
{\Huge\bfseries\color{ink}Ревью\par}
\vspace{3mm}
{\Huge\bfseries\color{ink}домашней работы\par}
\vspace{15mm}
{\large Ксения Клыкова}\par
\vspace{10mm}
{\normalsize\textcolor{muted}{31 октября 2026 г.}}\par
\vfill
{\normalsize Лёвин Артём Александрович}\par
\vspace{3mm}
{\small\textcolor{muted}{эксклюзивно для Ксении Клыковой}}\par
\vspace{11mm}
\begin{tikzpicture}
\draw[accent!55,line width=1.2pt] (0,0) .. controls (2.5,0.8) and (5.3,-0.8) .. (8.5,0.05) .. controls (10.2,0.55) and (12.1,0.36) .. (13.3,-0.15);
\end{tikzpicture}
\vspace{9mm}
\end{titlepage}
\pagenumbering{arabic}

\section*{Сводная таблица}
\textcolor{muted}{В таблице представлены только задания, требующие разбора. Номера открывают соответствующую страницу.}\par
\vspace{6mm}
{\small
\begin{tabular}{@{}p{10mm}p{27mm}p{54mm}p{22mm}p{23mm}@{}}
\toprule
\textbf{№} & \textbf{Тема} & \textbf{Суть затруднения} & \textbf{Ваш ответ} & \textbf{Верный ответ}\\
\midrule
\hyperref[task:8]{\textbf{8}} & Делители-квадраты & Не учтены одновременно чётность показателей и кратность 45 & 24 & \textbf{8}\\
\addlinespace[4pt]
\hyperref[task:9]{\textbf{9}} & Кратность и исключение & Решение среди присланных листов отсутствует & Не указан & \textbf{32}\\
\addlinespace[4pt]
\hyperref[task:10]{\textbf{10}} & Квадраты и ограничения & Решение среди присланных листов отсутствует & Не указан & \textbf{9}\\
\bottomrule
\end{tabular}\par
}
\vspace{8mm}
\textcolor{muted}{Примечание. По присланным фотографиям проверены задания № 1--8. Решений № 9 и № 10 на этих листах нет; они разобраны как невыполненные.}

\clearpage
\thispagestyle{empty}
\vspace*{0.34\textheight}
\begin{center}
{\Huge\bfseries\color{ink}Разбор ошибок}\par
\vspace{8mm}
{\large\textcolor{muted}{Делимость, показатели простых множителей и точный подсчёт}}\par
\end{center}

\clearpage
{\small\setlength{\parskip}{2pt}
\tasktitle{8}{Квадрат натурального числа, кратный 45}
\subhead{Текст задания}
Сколько делителей числа
\[
N=2^6\cdot3^5\cdot5^3
\]
одновременно кратны $45$ и являются квадратами натуральных чисел?

\subhead{Ваше решение}
Правильно записано $45=3^2\cdot5$ и представление делителя в виде
\[
d=2^x\cdot3^y\cdot5^z.
\]
Также найдены допустимые диапазоны для кратности $45$:
\[
0\le x\le6,\qquad 2\le y\le5,\qquad1\le z\le3.
\]
Затем число вариантов определено через произведение
\[
(3+1)(2+1)(1+1)=24.
\]
\textbf{Ваш ответ:} $24$.

\subhead{Где именно возникает ошибка}
\errorlabel{} Диапазоны показателей выписаны с учётом кратности $45$, однако на этапе подсчёта вариантов не соблюдено требование, чтобы делитель был \emph{квадратом натурального числа}. Для такого числа показатели всех простых множителей обязаны быть чётными. Кроме того, разрешённые чётные значения нужно считать внутри соответствующих диапазонов, начиная с их нижних границ.

\subhead{Почему это неверно}
Если $d=m^2$, то в разложении $d$ на простые множители каждый показатель имеет вид $2k$. Например, $2^2\cdot3^4$ является квадратом, а $2^2\cdot3^3$ --- нет. Поэтому количество всех целых показателей из интервала нельзя использовать вместо количества \emph{чётных} показателей. Условие кратности $45=3^2\cdot5$ дополнительно исключает значения $y<2$ и $z<1$.

\subhead{Ключевая идея и исправление}
\keylabel{} Для каждого простого множителя сначала найдите допустимый диапазон показателя, а затем оставьте только чётные показатели.
\[
\begin{array}{c|c|c}
\text{Показатель}&\text{Возможные чётные значения}&\text{Количество}\\\hline
x&0,2,4,6&4\\
y&2,4&2\\
z&2&1
\end{array}
\]
Правильное произведение: $4\cdot2\cdot1=8$.

\subhead{Полное правильное решение}
Любой положительный делитель $N$ имеет вид $2^x3^y5^z$ с целыми неотрицательными показателями $x\le6$, $y\le5$, $z\le3$. Из кратности $45=3^2\cdot5$ получаем $y\ge2$, $z\ge1$; на $x$ это условие не влияет. Так как $d$ --- квадрат, все три показателя чётные. Отсюда $x\in\{0,2,4,6\}$, $y\in\{2,4\}$, $z=2$. Выборы независимы, значит, всего $4\cdot2\cdot1=8$ делителей.

\subhead{Проверка}
Действительно, все они задаются формулой $d=2^{2a}\cdot3^{2b}\cdot5^2$, где $a\in\{0,1,2,3\}$, $b\in\{1,2\}$. Каждый такой делитель является квадратом, содержит множитель $3^2\cdot5$ и делит $N$. Других чётных показателей в заданных пределах нет.

\finallabel{} \textbf{Ответ: $8$.}

\subhead{Задача на закрепление}
Сколько делителей числа $2^5\cdot3^6\cdot5^4$ одновременно кратны $75$ и являются квадратами натуральных чисел?
}
\clearpage
\tasktitle{9}{Кратность 6, но отсутствие кратности 60}
\subhead{Текст задания}
Сколько делителей числа
\[
N=2^5\cdot3^4\cdot5^3
\]
делятся на $6$, но не делятся на $60$?

\subhead{Ваше решение и ответ}
На присланных листах решение задания № 9 отсутствует. \textbf{Ваш ответ:} не указан.

\subhead{Где именно возникает затруднение}
\errorlabel{} Отсутствует обоснование, которое одновременно учитывает оба требования: кратность $6$ и отсутствие кратности $60$. Последний правильный шаг и первый неверный переход указать невозможно, поскольку ход решения не представлен.

\subhead{Почему важно учитывать оба условия}
Если делитель кратен $6=2\cdot3$, в нём обязаны присутствовать хотя бы одна двойка и одна тройка. Однако часть таких делителей кратна и $60=2^2\cdot3\cdot5$ --- их нужно исключить. Простого подсчёта всех делителей, кратных $6$, недостаточно.

\subhead{Ключевая идея и исправление}
\keylabel{} Сначала посчитайте все делители, удовлетворяющие положительному условию, затем вычтите тех, которые удовлетворяют запрещённому условию.

\subhead{Полное правильное решение}
Любой делитель $N$ представим как
\[
d=2^a\cdot3^b\cdot5^c,\quad
0\le a\le5,\ 0\le b\le4,\ 0\le c\le3.
\]
\textbf{Шаг 1.} Для кратности $6$ необходимо $a\ge1$ и $b\ge1$. Показатель $a$ выбирается $5$ способами, $b$ --- $4$ способами, $c$ --- $4$ способами. Поэтому
\[
N_6=5\cdot4\cdot4=80.
\]
\textbf{Шаг 2.} Делитель кратен $60=2^2\cdot3\cdot5$, если $a\ge2$, $b\ge1$, $c\ge1$. Следовательно,
\[
N_{60}=4\cdot4\cdot3=48.
\]
Все делители, кратные $60$, уже входят в число $N_6$, поскольку число $60$ кратно $6$.

\textbf{Шаг 3.} Вычитаем запрещённые делители:
\[
N_6-N_{60}=80-48=32.
\]

\subhead{Проверка другим способом}
Чтобы делитель был кратен $6$, но не кратен $60$, возможны два непересекающихся случая: $a=1$ при любых $c=0,1,2,3$; либо $a\ge2$ и $c=0$. В обоих случаях $b=1,2,3,4$. Получаем
\[
1\cdot4\cdot4+4\cdot4\cdot1=16+16=32.
\]
Результаты совпадают.

\finallabel{} \textbf{Ответ: $32$.}

\subhead{Задача на закрепление}
Сколько делителей числа $2^6\cdot3^3\cdot5^2$ делятся на $6$, но не делятся на $30$?

\clearpage
\tasktitle{10}{Квадраты, кратные 72, но не кратные 64}
\subhead{Текст задания}
Сколько делителей числа
\[
N=2^7\cdot3^6\cdot5^4
\]
одновременно являются квадратами натуральных чисел, кратны $72$, но не кратны $64$?

\subhead{Ваше решение и ответ}
На присланных листах решение задания № 10 отсутствует. \textbf{Ваш ответ:} не указан.

\subhead{Где именно возникает затруднение}
\errorlabel{} Решение не представлено, поэтому нельзя выделить последний правильный и первый неправильный шаг. Для полного решения необходимо совместить три ограничения: чётность показателей, нижние границы кратности $72$ и верхнюю границу, вытекающую из отсутствия кратности $64$.

\subhead{Почему важно совместить ограничения}
Квадрат натурального числа имеет только чётные показатели в разложении на простые множители. Кратность $72=2^3\cdot3^2$ заставляет показатели двойки и тройки быть достаточно большими. Условие «не кратно $64=2^6$» исключает показатели двойки, равные $6$ и выше. Каждое ограничение обязательно.

\subhead{Ключевая идея и исправление}
\keylabel{} Выпишите допустимые показатели для каждого простого множителя как пересечение всех условий. После этого перемножьте количества вариантов.

\subhead{Полное правильное решение}
Пусть $d=2^a\cdot3^b\cdot5^c$ --- положительный делитель $N$. Тогда
\[
0\le a\le7,\qquad0\le b\le6,\qquad0\le c\le4.
\]
\textbf{Шаг 1.} Кратность $72=2^3\cdot3^2$ означает $a\ge3$, $b\ge2$.

\textbf{Шаг 2.} Некратность $64=2^6$ означает $a<6$.

\textbf{Шаг 3.} Для квадрата все показатели должны быть чётными. Значит,
\[
\begin{aligned}
a&\in\{4\}&&\text{(один вариант)},\\
b&\in\{2,4,6\}&&\text{(три варианта)},\\
c&\in\{0,2,4\}&&\text{(три варианта)}.
\end{aligned}
\]
Выборы независимы. Поэтому число подходящих делителей:
\[
1\cdot3\cdot3=9.
\]

\subhead{Проверка}
Каждый подходящий делитель имеет вид $2^4\cdot3^b\cdot5^c$, где $b=2,4,6$, $c=0,2,4$. При таких показателях делитель является квадратом; содержит множитель $2^3\cdot3^2$ и, следовательно, кратен $72$; содержит только четыре двойки, поэтому на $64$ не делится. Иных вариантов нет.

\finallabel{} \textbf{Ответ: $9$.}

\subhead{Задача на закрепление}
Сколько делителей числа $2^8\cdot3^4\cdot5^6$ являются квадратами натуральных чисел, кратны $18$, но не кратны $32$?

\end{document}